% ==========================================
% BAB II STUDI LITERATUR
% ==========================================
\cleardoublepage
\chapter{STUDI LITERATUR}
\label{chap:studi-literatur}

Literatur yang dikaji dalam bab ini memberikan landasan konseptual untuk
memahami domain inspeksi kontainer serta prinsip arsitektur sistem, sekaligus
menempatkan tugas akhir ini dalam konteks regulasi dan standar industri yang
berlaku.

\section{Sistem Inspeksi Kontainer: Konteks dan Tantangan}

Inspeksi kontainer merupakan elemen krusial dalam rantai logistik global yang
menjamin keamanan, kepatuhan regulasi, dan integritas barang dalam perdagangan
internasional. World Customs Organization (WCO) menempatkan inspeksi dan
pertukaran informasi kepabeanan sebagai bagian dari upaya mengamankan sekaligus
memfasilitasi perdagangan global \autocite{wco2025}. Dalam konteks Indonesia,
National Logistics Ecosystem (NLE) mengidentifikasi duplikasi dokumen dan
ketidakseimbangan informasi sebagai salah satu penyebab hambatan logistik,
sedangkan biaya logistik nasional masih menjadi isu efisiensi yang terus
diturunkan \autocite{kemenkeu2023nle,pelindo2023logistik}. Kondisi tersebut
membuat pencatatan inspeksi kontainer, keterhubungan data, dan ketersediaan
bukti pemeriksaan menjadi bagian penting dari perbaikan proses di pelabuhan.

Secara tradisional, inspeksi kontainer dilakukan melalui pemeriksaan manual oleh
petugas dengan dua kelompok acuan. Pertama, CTPAT 7-Point Container Inspection
Checklist dari U.S. Customs and Border Protection digunakan sebagai acuan
pemeriksaan keamanan dan integritas fisik kontainer sebelum pemuatan
\autocite{cbp2022}. Kedua, Guide for Container Equipment Inspection 6th Edition
dari Institute of International Container Lessors digunakan sebagai acuan
pemeriksaan kondisi dan kerusakan peralatan kontainer \autocite{iicl2016}.
Meskipun acuan tersebut membantu menstrukturkan pemeriksaan manual, pendekatan
manual tetap memiliki keterbatasan inheren: konsistensi rendah, tingginya
tingkat \textit{human error}, durasi waktu inspeksi yang lama, serta dokumentasi
yang sering tidak memiliki struktur metadata lengkap untuk ditelusuri kembali
\autocite{cbp2022,iicl2016,iso27037}.

Transformasi menuju pendekatan otomasi menuntut pemahaman terhadap teori
inspeksi kontainer, prinsip arsitektur sistem, dan kerangka tata kelola yang
mendukung operasi yang konsisten, terukur, dan berkelanjutan. Literatur yang ada
menyediakan fondasi teoretis untuk merancang sistem yang sesuai dengan kebutuhan
operasional dan regulasi pelabuhan.

\subsection{Klasifikasi Metode Inspeksi Kontainer}

World Customs Organization (WCO) mengklasifikasikan metode inspeksi kontainer ke dalam beberapa kategori berdasarkan tingkat intrusivitas dan prinsip dasar teknologi yang digunakan \autocite{wco2020}. Klasifikasi ini penting untuk memahami dampak operasional, biaya, dan kelayakan pendekatan inspeksi dalam konteks pelabuhan.

\begin{table}[ht]
  \centering
  \caption{Klasifikasi Metode Inspeksi Kontainer Menurut WCO}
  \label{tbl:inspeksi_klasifikasi}
  \begin{tabular}{p{0.25\textwidth} p{0.40\textwidth} p{0.25\textwidth}}
    \toprule
    \textbf{Kategori} & \textbf{Deskripsi} & \textbf{Karakteristik} \\
    \midrule
    \textbf{\textit{Intrusive Inspection}} & Pemeriksaan manual langsung pada isi kontainer dengan membuka segel dan memeriksa dokumen & Akurasi tinggi untuk validasi barang, tetapi memerlukan waktu signifikan dan tenaga kerja intensif \\
    \textbf{\textit{Non-Intrusive Inspection} (NII)} & Penggunaan teknologi pencitraan tanpa membuka kontainer untuk memeriksa isi dan struktur & Efisien waktu, dapat diotomatisasi, tetapi memerlukan investasi infrastruktur besar \\
    \textbf{\textit{Visual-Mechanical Inspection}} & Pemeriksaan struktur kontainer menggunakan panduan standar industri & Relatif cepat dan murah, tetapi bergantung pada subjektivitas dan pengalaman petugas \\
    \textbf{\textit{Sensor-Based Inspection}} & Penggunaan sensor untuk deteksi anomali berbasis parameter fisik & Dapat memberikan pemantauan waktu nyata, tetapi sensitif terhadap kondisi lingkungan \\
    \bottomrule
  \end{tabular}
\end{table}

\FloatBarrier

Tabel \ref{tbl:inspeksi_klasifikasi} menggambarkan empat pendekatan utama inspeksi kontainer yang umum digunakan dalam industri logistik. Setiap pendekatan memiliki karakteristik dan aplikasi yang berbeda, dan pemilihan pendekatan yang tepat bergantung pada kebutuhan operasional, regulasi, dan konteks pelabuhan.

\subsection{Model Alur Kerja Inspeksi Berbasis Risiko}

World Customs Organization SAFE \textit{Framework} 2025 menganjurkan pendekatan
inspeksi yang berbasis risiko dan \textit{data-driven} \autocite{wco2025}.
Model alur kerja berbasis risiko menggantikan pemeriksaan acak dengan model
yang terstruktur dan dapat diaudit. Tahap awalnya adalah \textit{targeting},
yaitu seleksi risiko berdasarkan profil data historis, karakteristik pengirim,
jenis barang, dan rute perdagangan untuk mengidentifikasi kontainer yang
memerlukan pemeriksaan lebih lanjut. Setelah itu, \textit{screening} dilakukan
melalui pemeriksaan awal menggunakan metode nonintrusif untuk kontainer
berisiko sedang sehingga proses dapat dipercepat tanpa pembongkaran fisik.
Hasil pemeriksaan kemudian masuk ke tahap \textit{analysis} untuk
mengidentifikasi anomali yang memerlukan tindak lanjut. Pada kontainer
berisiko tinggi, \textit{physical inspection} tetap diperlukan sebagai validasi
manual yang lebih komprehensif. Seluruh proses tersebut bermuara pada
\textit{compliance decision}, yaitu keputusan kepabeanan berdasarkan bukti
digital dan dokumentasi yang dapat diaudit.

Pola ini menekankan pendekatan berbasis risiko yang memerlukan sistem dengan
kemampuan untuk mengelola informasi secara terstruktur, menyediakan jejak audit
yang jelas, dan mendukung pengambilan keputusan yang konsisten. Setiap tahapan
dalam alur kerja menghasilkan informasi yang harus dikelola secara sistematis
untuk memastikan riwayat kontainer dan kepatuhan terhadap regulasi.

\section{Teknologi Inspeksi Kontainer Modern}

\textbf{Konsep \textit{Non-Intrusive Inspection} (NII) dan pencitraan.}
\textit{Non-Intrusive Inspection} (NII) merupakan konsep pemeriksaan isi
kontainer tanpa perlu membuka kontainer. Teknologi NII mencakup pendekatan
pencitraan yang dapat memetakan struktur dan isi kontainer secara
nondestruktif. Konsep dasar NII didasarkan pada prinsip fisika yang memungkinkan
penetrasi material untuk menghasilkan representasi visual dari interior objek
tanpa intervensi fisik langsung \autocite{wco2020}.

Pendekatan representasi visual digunakan untuk mengidentifikasi deformasi struktural pada permukaan kontainer melalui analisis pola. Literatur memaparkan konsep ini sebagai proses representasi visual yang digunakan untuk memahami pola dan karakteristik objek secara konsisten.

\textbf{Termografi inframerah sebagai metode nondestruktif.}
\textit{Infrared Thermography} (IRT) merupakan metode deteksi nondestruktif yang memanfaatkan perbedaan distribusi suhu pada permukaan objek untuk mengidentifikasi anomali. Prinsip dasar IRT didasarkan pada hukum Stefan-Boltzmann yang menyatakan bahwa semua objek dengan temperatur lebih tinggi daripada nol absolut memancarkan radiasi elektromagnetik \autocite{maldague2002}. Dalam konteks inspeksi kontainer, IRT dapat digunakan untuk mendeteksi anomali termal yang mengindikasikan kebocoran, deformasi struktural yang menyebabkan perubahan konduktivitas panas, atau gangguan internal yang mengakibatkan akumulasi panas \autocite{meola2004}.

Meskipun efektif sebagai \textit{screening} awal ketika terdapat perbedaan temperatur signifikan, IRT memiliki keterbatasan teoretis: kerusakan nontermal sering kali tidak terdeteksi karena tidak menghasilkan kontras temperatur yang memadai. Emisivitas permukaan, jarak pemindaian, dan kondisi lingkungan dapat memengaruhi akurasi hasil secara signifikan sehingga metode ini umumnya digunakan sebagai pelengkap metode inspeksi lainnya.

\section{Teori Arsitektur Sistem Terintegrasi}

\subsection{Prinsip Arsitektur Sistem Terdistribusi}

Arsitektur sistem terdistribusi untuk sistem otomasi inspeksi kontainer harus
mempertimbangkan prinsip-prinsip dasar dari \textit{software architecture} yang
telah dipaparkan oleh Bass, Clements, dan Kazman \autocite{bass2022}. Dalam
\textit{Software Architecture in Practice}, mereka menekankan bahwa arsitektur
bukan sekadar struktur komponen, tetapi juga mencakup prinsip desain yang
mendasari pembagian tanggung jawab sistem. Prinsip \textit{separation of
concerns} menuntut pemisahan tanggung jawab yang jelas antara komponen.
\textit{Modularity} mengarahkan sistem agar tersusun dari unit independen yang
dapat dikelola dan dikembangkan secara terpisah. \textit{Loose coupling}
menekan ketergantungan antarkomponen melalui antarmuka yang terdefinisi secara
eksplisit, sedangkan \textit{high cohesion} memastikan fungsi yang saling berkaitan
dikelompokkan dalam komponen yang sama.

Prinsip-prinsip ini membentuk fondasi untuk merancang sistem yang tidak hanya efektif secara fungsional tetapi juga dapat dipelihara, dikembangkan, dan diadaptasi sesuai dengan perubahan kebutuhan operasional.

\subsection{Atribut Kualitas dalam Arsitektur Sistem}

ISO/IEC 25010 mendefinisikan model kualitas untuk perangkat lunak yang
digunakan pada sistem otomasi inspeksi kontainer \autocite{iso25010}.
Dalam tugas akhir ini, atribut kualitas yang digunakan difokuskan pada
karakteristik yang paling langsung memengaruhi rancangan arsitektur sistem,
yaitu ketepatan fungsi inspeksi, kemudahan penggunaan, keamanan data, dan
kemudahan perubahan sistem.

\begin{table}[ht]
  \centering
  \small
  \caption{Atribut Kualitas ISO/IEC 25010 yang Digunakan dalam Perancangan}
  \label{tbl:quality_attributes}
  \begin{tabular}{p{0.24\textwidth} p{0.28\textwidth} p{0.38\textwidth}}
    \toprule
    \textbf{Atribut} & \textbf{Subatribut yang Digunakan} & \textbf{Implikasi terhadap Rancangan Arsitektur} \\
    \midrule
    \textbf{\textit{Functional suitability}} & \textit{Functional completeness}, \textit{functional correctness} & Arsitektur harus memungkinkan seluruh hasil inspeksi utama dicatat dan dikaitkan dengan entitas kontainer yang benar \\
    \textbf{\textit{Usability}} & \textit{Operability}, \textit{learnability} & Antarmuka dan alur informasi harus mudah dioperasikan oleh pengguna operasional tanpa menambah kompleksitas pemeriksaan \\
    \textbf{\textit{Security}} & \textit{Integrity}, \textit{authenticity} & Data hasil inspeksi, bukti visual, dan keputusan pemeriksaan harus terlindungi dari perubahan tidak sah dan memiliki asal data yang jelas \\
    \textbf{\textit{Maintainability}} & \textit{Modifiability}, \textit{testability} & Komponen sistem perlu dipisahkan secara modular agar perubahan pada \textit{edge}, API, penyimpanan, atau antarmuka dapat diuji dan dikembangkan secara terkendali \\
    \bottomrule
  \end{tabular}
\end{table}

\FloatBarrier

Atribut kualitas pada Tabel \ref{tbl:quality_attributes} digunakan
sebagai landasan untuk mendefinisikan kebutuhan nonfungsional sistem inspeksi
kontainer. Hubungan antara atribut utama dan subatribut yang digunakan dalam
perancangan ditunjukkan pada Gambar \ref{fig:iso25010_quality_model}.

\begin{figure}[htbp]
  \centering
  \resizebox{0.90\textwidth}{!}{%
  \begin{tikzpicture}[
      font=\small,
      core/.style={rectangle, draw=black, rounded corners, align=center, minimum width=4.1cm, minimum height=1.1cm},
      attr/.style={rectangle, draw=black, rounded corners, align=center, minimum width=3.6cm, minimum height=0.95cm},
      detail/.style={rectangle, draw=black, align=center, minimum width=3.6cm, minimum height=0.82cm},
      flow/.style={-{Latex[length=2mm]}, thick},
      bus/.style={thick}
    ]
    \node[core] (system) at (0,0) {\textbf{Kualitas Arsitektur}\\Sistem Inspeksi Kontainer};

    \node[attr] (functional) at (-6.0,-1.95) {\textit{Functional}\\\textit{Suitability}};
    \node[attr] (usability) at (-2.0,-1.95) {\textit{Usability}};
    \node[attr] (security) at (2.0,-1.95) {\textit{Security}};
    \node[attr] (maintainability) at (6.0,-1.95) {\textit{Maintainability}};

    \node[detail] (functionalDetail) at (-6.0,-3.60) {Hasil inspeksi\\tercatat lengkap};
    \node[detail] (usabilityDetail) at (-2.0,-3.60) {Alur kerja dapat\\dioperasikan petugas};
    \node[detail] (securityDetail) at (2.0,-3.60) {Data dan bukti\\terlindungi};
    \node[detail] (maintainabilityDetail) at (6.0,-3.60) {Komponen dapat\\diubah bertahap};

    \node[detail, minimum width=7.9cm, minimum height=0.88cm] (evaluation) at (0,-5.55)
      {Digunakan sebagai dasar kebutuhan nonfungsional\\dan kriteria evaluasi rancangan Cargovision};

    \draw[bus] (system.south) -- (0,-1.15);
    \draw[bus] (-6.0,-1.15) -- (6.0,-1.15);
    \draw[flow] (-6.0,-1.15) -- (functional.north);
    \draw[flow] (-2.0,-1.15) -- (usability.north);
    \draw[flow] (2.0,-1.15) -- (security.north);
    \draw[flow] (6.0,-1.15) -- (maintainability.north);
    \draw[flow] (functional.south) -- (functionalDetail.north);
    \draw[flow] (usability.south) -- (usabilityDetail.north);
    \draw[flow] (security.south) -- (securityDetail.north);
    \draw[flow] (maintainability.south) -- (maintainabilityDetail.north);
    \draw[bus] (functionalDetail.south) -- (-6.0,-4.55);
    \draw[bus] (usabilityDetail.south) -- (-2.0,-4.55);
    \draw[bus] (securityDetail.south) -- (2.0,-4.55);
    \draw[bus] (maintainabilityDetail.south) -- (6.0,-4.55);
    \draw[bus] (-6.0,-4.55) -- (6.0,-4.55);
    \draw[flow] (0,-4.55) -- (evaluation.north);
  \end{tikzpicture}}
  \caption{Model kualitas ISO/IEC 25010 yang digunakan sebagai landasan untuk mendefinisikan kebutuhan nonfungsional sistem pemeriksaan kontainer}
  \label{fig:iso25010_quality_model}
\end{figure}
\FloatBarrier

Atribut kualitas ini membentuk kriteria evaluasi untuk desain
arsitektur dan menjadi dasar untuk pengambilan keputusan desain sepanjang
proses perancangan sistem.

\subsection{Paradigma \textit{Edge-Cloud Computing}}

Konsep \textit{edge-cloud computing} mengacu pada model arsitektur komputasi yang membagi beban pemrosesan antara komponen di lokasi sumber data dan pusat pemrosesan terpusat \autocite{shi2016}. Shi dkk. mendefinisikan \textit{edge computing} sebagai paradigma yang menempatkan komputasi dan penyimpanan data di dekat sumber data untuk mengoptimalkan respons dan efisiensi sistem.

Paradigma \textit{edge-cloud} menekankan beberapa prinsip utama. Prinsip
\textit{locality} menempatkan pemrosesan data di dekat sumber untuk
mengoptimalkan penggunaan \textit{bandwidth} dan mengurangi latensi.
\textit{Autonomy} menuntut komponen tepi tetap dapat beroperasi secara
terbatas ketika koneksi dengan pusat terganggu. \textit{Scalability}
dipenuhi melalui pusat pemrosesan yang menyediakan kapabilitas elastis untuk
analisis komprehensif dan penyimpanan jangka panjang. Sementara itu,
\textit{efficiency} dicapai melalui pembagian beban pemrosesan antara
komponen lokal dan terpusat.

Varghese dkk. memperluas konsep ini dengan mengidentifikasi pola-pola distribusi beban kerja untuk aplikasi dengan kebutuhan respons cepat pada pengolahan data terdistribusi \autocite{varghese2016}. Secara teoretis, konsep ini menekankan pentingnya pembagian fungsional antara komponen lokal dan terpusat sesuai karakter beban.

\subsection{Prinsip Arsitektur Layanan}

Newman mendefinisikan arsitektur layanan sebagai pendekatan pengembangan yang
mengorganisir aplikasi menjadi koleksi layanan yang secara fungsional
terdekomposisi, \textit{loosely coupled}, dan dapat dikelola secara independen
\autocite{newman2015}. Prinsip \textit{service autonomy} memberikan kendali
kepada setiap layanan terhadap logika dan data yang dikelolanya. Prinsip
\textit{service boundaries} menuntut batas layanan yang jelas berdasarkan
domain fungsional atau bisnis. Komunikasi antarlayanan dijaga melalui
\textit{contract-based communication}, yaitu antarmuka yang terdefinisi secara
formal. Di sisi tata kelola, \textit{decentralized governance} memungkinkan
setiap layanan menggunakan pendekatan yang sesuai dengan kebutuhan domainnya.

Pendekatan ini memungkinkan sistem untuk berkembang secara evolusioner karena komponen dapat dimodifikasi atau diganti tanpa memerlukan perubahan menyeluruh pada sistem. Prinsip ini digunakan dalam literatur untuk menjelaskan pola adaptasi arsitektur terhadap perubahan regulasi dan kebutuhan operasional.

\section{Kecerdasan Buatan dalam Inspeksi Visual}

\textbf{Konsep pemodelan visual otomatis.}
\textit{Computer vision} merupakan bidang yang mempelajari bagaimana sistem komputasi dapat memperoleh pemahaman tingkat tinggi dari citra atau video digital \autocite{szeliski2022}. Prinsip dasar melibatkan transformasi informasi visual mentah menjadi representasi semantik yang dapat digunakan untuk pengambilan keputusan. Konsep ini dipaparkan dalam literatur sebagai pendekatan untuk memahami karakteristik interpretasi visual yang konsisten. Dalam konteks inspeksi kontainer, pendekatan ini digunakan karena kerusakan struktural memiliki pola tepi, bentuk, dan tekstur yang dapat dipelajari model secara sistematis.

\textbf{Pembelajaran dari pola data visual.}
Pendekatan pembelajaran mesin untuk deteksi objek didasarkan pada prinsip bahwa model dapat membangun representasi konseptual dari pola visual untuk membedakan karakteristik objek dalam berbagai kondisi \autocite{goodfellow2016}. Representasi ini memungkinkan sistem untuk mengenali kesamaan dan perbedaan dalam informasi visual tanpa memerlukan spesifikasi eksplisit untuk setiap kemungkinan variasi.

Literatur membedakan pendekatan deteksi berdasarkan organisasi proses analisis visual, yang dapat dilakukan secara bertahap dengan pemisahan tahap analisis awal dari tahap identifikasi objek, atau secara terintegrasi dengan penggabungan kedua proses. Literatur menyoroti adanya pertimbangan antara kedalaman analisis dan efisiensi proses dalam pemilihan pendekatan deteksi visual.

\textbf{Prinsip keandalan model prediktif.}
Validasi model pembelajaran mesin merupakan proses memastikan bahwa model dapat melakukan generalisasi terhadap data yang belum pernah dilihat sebelumnya. Keandalan model diukur melalui prediksi yang konsisten dan akurat dalam kondisi operasional yang bervariasi. Proses validasi melibatkan evaluasi sistematis terhadap kinerja model menggunakan data yang tidak identik dengan data pembentuk representasi model, serta analisis terhadap pola kesalahan untuk mengidentifikasi area perbaikan.

Literatur menekankan bahwa keandalan model merupakan aspek penting dalam kajian pembelajaran mesin, terutama pada domain yang memerlukan konsistensi interpretasi visual.

\section{Komunikasi Data dan Integrasi Sistem}

\textbf{Konsep komunikasi berbasis kejadian.}
Arsitektur berbasis kejadian merupakan pola desain yang membuat komponen sistem berinteraksi melalui notifikasi perubahan status, memungkinkan independensi temporal dan fungsional antara komponen \autocite{hohpe2003}. Dalam pendekatan ini, komponen yang menghasilkan informasi tidak perlu mengetahui komponen mana yang akan merespons informasi tersebut, dan sebaliknya. Konsep ini menciptakan sistem yang lebih fleksibel karena komponen baru dapat ditambahkan untuk merespons kejadian yang sudah ada tanpa mengubah komponen penghasil kejadian. Pendekatan berbasis kejadian ini selaras dengan kebutuhan inspeksi kontainer yang memerlukan penelusuran perubahan status kontainer sepanjang alur pemeriksaan secara transparan.

Pendekatan berbasis kejadian memungkinkan sistem untuk berkembang secara evolusioner dengan menambahkan respons baru terhadap kejadian yang ada, mendukung adaptasi sistem terhadap kebutuhan yang berubah.

\textbf{Mekanisme penyampaian informasi terdistribusi.}
Sistem terdistribusi memerlukan mekanisme untuk menyampaikan informasi antarkomponen yang mungkin beroperasi pada waktu yang berbeda atau mengalami gangguan komunikasi \autocite{kleppmann2017}. Konsep penyampaian informasi yang toleran terhadap gangguan melibatkan penggunaan perantara yang mengatur pengaliran informasi antarkomponen. Pendekatan ini meningkatkan ketahanan sistem terhadap kegagalan sementara pada komponen individual.

Mekanisme ini memungkinkan komponen untuk beroperasi secara independen tanpa bergantung pada ketersediaan simultan semua komponen yang terlibat, mendukung operasi yang berkelanjutan meskipun terjadi gangguan pada bagian tertentu dari sistem.

\textbf{Prinsip interoperabilitas dan standardisasi.}
Interoperabilitas sistem terdistribusi memerlukan kontrak formal yang
memungkinkan komponen heterogen untuk berkomunikasi secara efektif
\autocite{bass2022}. Bass dkk. menekankan bahwa standardisasi antarmuka
merupakan kunci untuk mencapai interoperabilitas yang berkelanjutan.

Standardisasi tersebut mencakup \textit{interface contracts} yang mendefinisikan
operasi, parameter, dan semantik komunikasi antarkomponen; \textit{data format
standards} yang menjaga konsistensi interpretasi data; \textit{protocol
standards} yang memastikan keselarasan mekanisme pertukaran informasi; serta
\textit{versioning strategy} yang mengelola evolusi antarmuka tanpa merusak
keselarasan dengan komponen yang sudah ada.

Pendekatan kontrak standar memungkinkan komponen untuk berkembang selama
tetap mematuhi antarmuka yang telah disepakati. Prinsip ini mendukung evolusi
sistem jangka panjang tanpa menuntut perubahan serentak pada seluruh komponen.

\section{Tata Kelola Teknologi Informasi untuk Sistem Inspeksi}

\subsection{Kerangka COBIT untuk Tata Kelola TI}

COBIT (Control Objectives for Information and Related Technologies) merupakan kerangka kerja tata kelola dan manajemen teknologi informasi yang dikembangkan oleh ISACA \autocite{isaca2019}. COBIT 2019 menyediakan prinsip-prinsip untuk menyelaraskan TI dengan tujuan bisnis organisasi dan memastikan penggunaan teknologi yang efektif dan bertanggung jawab.

Prinsip dasar COBIT menekankan bahwa sistem harus dirancang untuk memenuhi
kebutuhan pemangku kepentingan yang beragam, mencakup organisasi dan ekosistem
terkait secara menyeluruh, serta mengintegrasikan tata kelola TI dengan tata
kelola organisasi. COBIT juga menekankan pendekatan holistik terhadap faktor
yang memengaruhi tata kelola dan pemisahan yang jelas antara fungsi tata kelola
dan fungsi manajemen operasional.

Dalam konteks sistem inspeksi kontainer, COBIT menyediakan kerangka untuk memastikan bahwa sistem tidak hanya efektif secara teknis tetapi juga selaras dengan tujuan operasional pelabuhan dan memenuhi persyaratan regulasi.

\textbf{ITIL untuk \textit{service management}.}
ITIL (Information Technology Infrastructure Library) merupakan kerangka kerja praktik terbaik untuk manajemen layanan TI yang dikembangkan oleh AXELOS \autocite{axelos2019}. ITIL 4 memperkenalkan konsep \textit{Service Value System} yang menekankan penciptaan nilai melalui layanan TI.

Komponen utama ITIL 4 mencakup \textit{Service Value Chain} sebagai model
operasi untuk merespons permintaan dan menciptakan nilai, \textit{Guiding
Principles} sebagai panduan umum dalam manajemen layanan, \textit{Practices}
sebagai kumpulan kemampuan organisasi, serta \textit{Continual Improvement}
sebagai pendekatan iteratif untuk meningkatkan layanan, komponen, dan aktivitas
secara berkelanjutan.

ITIL menyediakan kerangka untuk mengelola siklus hidup layanan sistem inspeksi kontainer, dari perencanaan hingga operasi dan peningkatan berkelanjutan.

\textbf{Prinsip \textit{audit trail} dan \textit{digital evidence}.}
ISO/IEC 27037:2012 menetapkan pedoman untuk identifikasi, pengumpulan, akuisisi, dan preservasi bukti digital \autocite{iso27037}. Dalam konteks sistem inspeksi kontainer, prinsip-prinsip \textit{audit trail} yang kuat sangat penting untuk memastikan integritas dan keabsahan data inspeksi.

Prinsip dasar \textit{audit trail} mencakup \textit{chain of custody} sebagai
dokumentasi kronologis atas kendali dan perpindahan bukti digital,
\textit{integrity verification} untuk memastikan data tidak dimodifikasi secara
tidak sah, \textit{non-repudiation} untuk membuktikan asal dan integritas data,
\textit{immutability} untuk menjaga catatan audit tidak berubah setelah
dicatat, serta \textit{accessibility} agar jejak audit dapat diakses untuk
keperluan investigasi atau kepatuhan.

Prinsip-prinsip ini membentuk dasar teoretis untuk mekanisme pencatatan log dan dokumentasi yang dapat dipertanggungjawabkan secara hukum dan operasional.

\section{Standar Industri dan Regulasi}

\textbf{Standar inspeksi kontainer internasional.}
Institute of International Container Lessors (IICL) dan \textit{Unified Container Inspection and Repair Criteria} (UCIRC) menetapkan pedoman klasifikasi kerusakan dan dokumentasi yang menjadi dasar konsistensi penilaian kontainer \autocite{iicl2016, icswsc2023}. Standar ini mendefinisikan pola penilaian kerusakan yang konsisten, kategori kerusakan berdasarkan jenis dan tingkat keparahan, serta pentingnya dokumentasi yang mengikuti standar untuk keperluan audit dan klaim.

\textbf{\textit{WCO SAFE Framework}.}
World Customs Organization SAFE \textit{Framework of Standards} merupakan kerangka kerja internasional untuk mengamankan dan memfasilitasi perdagangan global \autocite{wco2025}. Kerangka ini menyediakan prinsip-prinsip untuk inspeksi kontainer berbasis risiko dan pertukaran informasi elektronik antara otoritas kepabeanan.

Pilar utama WCO SAFE \textit{Framework} mencakup pertukaran informasi intelijen dan
risiko antarotoritas kepabeanan melalui \textit{Customs-to-Customs Network},
kolaborasi dengan sektor swasta melalui \textit{Customs-to-Business
Partnership}, pendekatan sistematis untuk mengelola risiko melalui \textit{Risk
Management}, serta penggunaan teknologi modern untuk inspeksi dan verifikasi
kontainer melalui pilar \textit{Technology and Equipment}.

Kerangka ini menyediakan konteks regulasi yang harus dipertimbangkan dalam merancang sistem inspeksi kontainer yang sesuai dengan standar internasional.

\textbf{Regulasi nasional Indonesia.}
Peraturan Menteri Perhubungan Nomor PM~8 Tahun 2022 mengatur tata cara
pelayanan kapal melalui Inaportnet sebagai sistem layanan elektronik di
pelabuhan \autocite{kemenhub2022}. Regulasi ini menunjukkan bahwa layanan
operasional pelabuhan diarahkan untuk menggunakan sistem informasi elektronik
sebagai kanal pemrosesan dan pertukaran data layanan kapal.

Dalam konteks kepabeanan, PER-11/BC/2020 mengatur perubahan atas tata cara
penyerahan, penatausahaan, perbaikan, dan pembatalan pemberitahuan rencana
kedatangan sarana pengangkut serta manifes kedatangan dan keberangkatan sarana
pengangkut \autocite{bc2020}. Regulasi-regulasi ini menjadi konteks bahwa
rancangan sistem otomasi inspeksi kontainer perlu memperhatikan keterhubungan
dengan dokumen dan sistem elektronik yang sudah digunakan dalam proses
pelabuhan.

\section{Prinsip Pengelolaan Data dalam Sistem Inspeksi}

Pengelolaan data merupakan aspek fundamental dalam sistem otomasi inspeksi kontainer yang menjamin keandalan informasi sepanjang siklus operasional. Literatur menunjukkan bahwa kualitas data menjadi faktor kritis dalam mendukung pengambilan keputusan yang akurat dan dapat dipertanggungjawabkan dalam konteks inspeksi kontainer.

Konsep kualitas data mencakup beberapa dimensi utama yang digunakan dalam
sistem inspeksi. \textit{Accuracy} mengacu pada tingkat kesesuaian antara data
yang tercatat dengan kondisi aktual yang direpresentasikan.
\textit{Completeness} merujuk pada kelengkapan atribut data yang diperlukan
untuk mendukung proses inspeksi dan audit. \textit{Consistency} menekankan
keseragaman representasi data di seluruh komponen sistem yang berbeda. Dimensi
kualitas data ini penting dalam literatur arsitektur karena menentukan
efektivitas alur informasi yang menjadi dasar bagi operasi inspeksi.

Prinsip \textit{metadata} menjadi penting dalam konteks sistem terdistribusi
ketika informasi tentang data itu sendiri diperlukan untuk memastikan
interpretasi yang tepat. \textit{Metadata} mencakup informasi tentang sumber
data, waktu akuisisi, metode pengumpulan, dan transformasi yang telah
dilakukan. Dalam alur kerja inspeksi, \textit{metadata} memungkinkan pelacakan
asal-usul informasi dan memfasilitasi audit terhadap proses pengambilan
keputusan.

Dalam tugas akhir ini, \textit{traceability} dipahami sebagai riwayat kontainer,
yaitu kemampuan untuk menelusuri hubungan antara identitas kontainer, waktu
akuisisi data, hasil inspeksi, bukti visual, dan perubahan status dari awal
pemeriksaan sampai informasi digunakan. Riwayat kontainer memungkinkan
verifikasi terhadap validitas hasil inspeksi dan mendukung akuntabilitas dalam
proses kepabeanan.

\section{Celah Literatur dan Kontribusi Penelitian}

\subsection{Identifikasi Celah Penelitian}

Review sistematis terhadap literatur yang ada mengidentifikasi beberapa celah
kajian yang signifikan:

\begin{enumerate}
  \item Kerangka arsitektur konseptual yang menyatukan komponen inspeksi masih
  belum banyak
  dibahas. Meskipun literatur tentang teknologi inspeksi individual, seperti
  NII, IRT, dan \textit{computer vision}, telah berkembang, belum ada kerangka
  arsitektur konseptual yang mengintegrasikan seluruh komponen dalam alur
  informasi kontainer.
  \item Prinsip pembagian fungsional \textit{edge-cloud} untuk domain inspeksi
  kontainer belum banyak dibahas pada konteks pelabuhan. Literatur
  \textit{edge-cloud computing} telah mapan secara teoretis, tetapi belum banyak mendefinisikan
  pembagian fungsional optimal antara pemrosesan lokal dan terpusat untuk
  \textit{workload} inspeksi kontainer.
  \item Model atribut kualitas untuk sistem inspeksi kontainer belum dirumuskan
  secara spesifik. ISO/IEC 25010 menyediakan kerangka umum, tetapi belum
  menjelaskan prioritas dan \textit{trade-off} atribut kualitas dalam konteks
  sistem otomasi inspeksi kontainer.
  \item Prinsip \textit{audit trail} multi-pemangku kepentingan belum
  teradaptasi secara spesifik untuk alur kerja inspeksi kontainer. ISO 27037
  menyediakan pedoman penanganan bukti digital secara umum, tetapi belum
  menjelaskan kebutuhan pelabuhan, bea cukai, pemilik barang, dan regulator
  secara simultan.
  \item Kerangka interoperabilitas untuk ekosistem pelabuhan Indonesia belum
  terdefinisi secara spesifik, terutama untuk standar antarmuka dan protokol
  komunikasi yang menghubungkan sistem inspeksi dengan Inaportnet, TOS, dan PCS.
\end{enumerate}

\subsection{Kontribusi Tugas Akhir}

Tugas akhir ini berupaya mengisi celah tersebut melalui lima kontribusi utama:

\begin{enumerate}
  \item Menyusun \textit{blueprint} arsitektur konseptual terintegrasi yang
  menggambarkan struktur komponen, alur informasi, dan prinsip integrasi untuk
  sistem otomasi inspeksi kontainer.
  \item Merumuskan prinsip pembagian fungsi \textit{edge-cloud} agar pemrosesan
  lokal dan terpusat dapat ditempatkan sesuai kebutuhan responsivitas, efisiensi,
  dan akurasi.
  \item Mengadaptasi model atribut kualitas ISO/IEC 25010 ke dalam konteks
  inspeksi kontainer sehingga prioritas dan \textit{trade-off} antaratribut dapat
  dijelaskan terhadap operasi pelabuhan.
  \item Merumuskan kerangka \textit{audit trail} dan \textit{chain of custody}
  digital untuk menjaga riwayat kontainer dan bukti inspeksi dalam lingkungan yang
  melibatkan banyak pemangku kepentingan.
  \item Mengusulkan kerangka interoperabilitas kontekstual yang menekankan
  standardisasi antarmuka dan protokol komunikasi sesuai regulasi nasional dan
  karakteristik ekosistem pelabuhan Indonesia.
\end{enumerate}

Identifikasi celah literatur dan kontribusi tersebut memperkuat perlunya
analisis kebutuhan sistem untuk merancang arsitektur yang sesuai dengan konteks
pelabuhan Indonesia. Bab III akan menguraikan analisis masalah, kebutuhan
fungsional dan nonfungsional, serta spesifikasi konseptual yang menjadi dasar
perancangan arsitektur sistem otomasi inspeksi kontainer.
